\documentclass[10pt]{article}
\usepackage[T1]{fontenc}
\usepackage[utf8]{inputenc}
\usepackage[french]{babel}
\usepackage{amsmath,amssymb,mathrsfs,mathpazo}
\usepackage{geometry}
\geometry{a4paper,margin=2cm}
\usepackage{graphicx,array,tabularx,tabularray,colortbl,multirow,multicol,booktabs,enumitem,pifont,ulem,fancybox,diagbox,cancel}
\usepackage{pstricks,pstricks-add,pst-plot,pst-tree,pst-node,pst-text,pst-func,pst-math,pst-bspline,pst-3dplot,pst-eucl,pst-solides3d}
\let\pscalculate\undefined
\usepackage{pst-bezier}
\usepackage{tikz,pgfplots,tkz-tab}
\usetikzlibrary{arrows,arrows.meta,patterns,calc,positioning,trees,decorations.markings,shapes.geometric,shapes.misc,shapes.arrows,intersections,decorations.pathmorphing,decorations.pathreplacing,backgrounds,shadows,babel}
\pgfplotsset{compat=1.18}
\usepgfplotslibrary{fillbetween}
\usepackage[np]{numprint}
\usepackage{listings,makecell,pas-tableur,twoopt}
\let\barre\undefined
\usepackage[pstricks]{tablvar}
\definecolor{darkgreen}{RGB}{0,100,0}
\providecommand{\red}{\color{red}}
\providecommand{\blue}{\color{blue}}
\providecommand{\green}{\color{green}}
\providecommand{\black}{\color{black}}
\providecommand{\white}{\color{white}}
\providecommand{\gray}{\color{gray}}
\providecommand{\cyan}{\color{cyan}}
\providecommand{\magenta}{\color{magenta}}
\providecommand{\yellow}{\color{yellow}}
\providecommand{\violet}{\color{violet}}
\providecommand{\orange}{\color{orange}}
\providecommand{\euro}{\texteuro}
\providecommand{\EUR}{\texteuro}
\providecommand{\ieme}{\textsuperscript{e}}
\providecommand{\degres}{\ensuremath{^\circ}}
\providecommand{\degree}{\ensuremath{^\circ}}
\providecommand{\vv}[1]{\overrightarrow{#1}}
\providecommand{\Coord}[1]{#1}
\providecommand{\decosix}{}
\providecommand{\starredbullet}{\ensuremath{\bullet}}
\pagestyle{empty}
\setlength{\parindent}{0pt}
\usepackage{adjustbox}\DeclareUnicodeCharacter{2236}{:}\DeclareUnicodeCharacter{2190}{\ensuremath{\leftarrow}}
\providecommand{\R}{\mathbb{R}}
\providecommand{\N}{\mathbb{N}}
\providecommand{\Z}{\mathbb{Z}}
\providecommand{\Q}{\mathbb{Q}}
\providecommand{\D}{\mathbb{D}}
\providecommand{\C}{\mathbb{C}}
\providecommand{\e}{{\mathrm{e}}}
\providecommand{\vect}[1]{\overrightarrow{#1}}
\providecommand{\vectt}[1]{\overrightarrow{\mathrm{#1}}}
\providecommand{\vvt}[1]{\overrightarrow{\mathrm{#1}}}
\providecommand{\np}[2][]{{#2~#1}}
\providecommand{\vv}[1]{{\overrightarrow{#1}}}
\renewcommand{\d}{{\mathrm{d}}}
\providecommand{\degres}{{^\circ}}
\providecommand{\degre}{{^\circ}}
\providecommand{\blue}{\color{blue}}
\providecommand{\red}{\color{red}}
\providecommand{\enskip}{\quad}
\providecommand{\llbracket}{\lbrack\!\lbrack}
\providecommand{\rrbracket}{\rbrack\!\rbrack}
\providecommand{\euro}{{\text{€}}}
\providecommand{\no}{{\text{n°}}}
\providecommand{\interval}[2]{{[#1;#2]}}
\providecommand{\intervFO}[2]{{[#1;#2[}}
\providecommand{\dotfill}{\ldots}

\providecommand{\ds}{\displaystyle}
\providecommand{\textding}[1]{#1}
\providecommand{\Oij}{\(\mathrm{O};\vec{\imath},\vec{\jmath}\)}
\providecommand{\Oijk}{\(\mathrm{O};\vec{\imath},\vec{\jmath},\vec{k}\)}
\providecommand{\Ouv}{\(\mathrm{O};\vec{u},\vec{v}\)}
\providecommand{\intervFF}[2]{\ensuremath{[#1;#2]}}
\renewcommand{\euro}{\ifmmode\text{\texteuro}\else\texteuro\fi}
\DeclareUnicodeCharacter{00B0}{\ifmmode^{\circ}\else\textdegree\fi}
\DeclareUnicodeCharacter{20AC}{\ifmmode\text{\texteuro}\else\texteuro\fi}
% Énoncés : Éducation nationale. Transcription et corrigés : APMEP, Denis Vergès et contributeurs. Découpage et classement : CoopMaths, Sébastien Lozano, Éric Elter et contributeurs. Adaptation de la présentation : Maths Explorer.
% CC BY-SA : https://coopmaths.fr/www/about/
% https://coopmaths.fr/alea/static/bac/2022/tex/bac_2022_05_sujet1_polynesie_2_cor.tex
\begin{document}
\noindent\textit{APMEP / CoopMaths. CC BY-SA. Présentation adaptée par Maths Explorer.}\par\medskip
\medskip

Selon les autorités sanitaires d'un pays, 7\,\% des habitants sont affectés par une certaine maladie.

Dans ce pays, un test est mis au point pour détecter cette maladie. Ce test a les caractéristiques suivantes :
\setlength\parindent{1cm}
\begin{itemize}
\item[$\bullet~~$] Pour les individus malades, le test donne un résultat négatif dans $20 \,\%$ des cas ;
\item[$\bullet~~$] Pour les individus sains, le test donne un résultat positif dans $1\,\%$ des cas.
\end{itemize}
\setlength\parindent{0cm}

Une personne est choisie au hasard dans la population et testée.

On considère les évènements suivants :

\setlength\parindent{1cm}
\begin{itemize}
\item[$\bullet~~$] $M$ \og la personne est malade \fg{} ;
\item[$\bullet~~$]  $T$ \og le test est positif \fg{}.
\end{itemize}
\setlength\parindent{0cm}

\medskip

\begin{enumerate}
\item Pour calculer la probabilité de l'évènement $M \cap T$, on s'appuie sur un arbre pondéré :



\begin{center}
\pstree[treemode=R,nodesepA=0pt,nodesepB=2.5pt,treesep=1cm,levelsep=2.5cm]{\TR{}}
{\pstree{\TR{$M$~}\taput{$0,07$}}
	{\TR{$T$}\taput{$0,8$}
	\TR{$\overline{T}$}\tbput{$0,2$}
	}
\pstree{\TR{$\overline{M}$~}\tbput{$0,93$}}
	{\TR{$T$}\taput{$0,01$}
	\TR{$\overline{T}$}\tbput{$0,99$}
	}
}
\end{center}



On a donc $P(M\cap T) = 0,07 \times 0,8 = 0,056$

\item 
$M$ et $\overline{M}$ formant une partition de l'univers, d'après les probabilités totales,

$P(T)=P(M \cap T)+ P\left(\overline{M} \cap T\right) = 0,07 \times 0,8 + 0,93 \times 0,01 = \np{0,0653}$.

\item 
$P_M(T)$ est la probabilité d'avoir un test positif sachant qu'on est malade, et  $P_T(M)$ est la probabilité d'être malade sachant que le test est positif.

Dans un contexte de dépistage de la maladie, il est plus pertinent de calculer la probabilité d'être malade sachant que le test est positif c'est-à-dire $P_T(M)$.

\item On considère dans cette question que la personne choisie au hasard a eu un test positif.

La probabilité qu'elle soit malade est:
$P_T(M)=\dfrac{P(M \cap T)}{P(T)}=\dfrac{0,07 \times 0,8}{0,0653} \approx 0,86$

\item On choisit des personnes au hasard dans la population. La taille de la population de ce pays permet d'assimiler ce prélèvement à un tirage avec remise.

On note $X$ la variable aléatoire qui donne le nombre d'individus ayant un test positif parmi les 10 personnes. 

	\begin{enumerate}
		\item 
Le tirage étant assimilé à un tirage avec remise, on a une indépendance des 10  tirages, le succès étant défini par le test positif, on peut assimiler cette variable aléatoire à une loi binomiale de paramètre $n=10$ et $p= \np{0,0653}$

		\item La probabilité pour qu'exactement deux personnes aient un test positif est:
		
$P(X = 2) = \binom{10}{2} \times 0,0653^2 \times (1-0,0653)^8 \approx 0,11$
	\end{enumerate}	

\item On veut déterminer le nombre minimum de personnes à tester dans ce pays pour que la probabilité qu'au moins l'une d'entre elles ait un test positif, soit supérieure à $99\,\%$.

Soit $n$ le nombre de personnes testées;  on cherche $n$ pour que $P(X\geqslant 1) \geqslant 0,99$.

On a $P(X\geqslant 1)=1 - P(X = 0)$  soit 

$1 - 0,9347^n \geqslant 0,99 \iff 0,9347^n \leqslant 0,01 \iff n \ln (0,9347) \leqslant \ln(0,01) \iff n \geqslant \dfrac{\ln(0,01}{\ln(0,9347} $. Or $\dfrac{\ln(0,01}{\ln(0,9347} \approx 68,2$.

Il faut donc tester $69$ personnes dans ce pays pour que au moins une ait un test positif.

\end{enumerate}

\vspace{1cm}

\end{document}
